%% ===================================================================
\section{States, matrix elements and Fourier transforms}\label{app:conv}
%% ===================================================================

\paragraph*{States.} The soft vacuum $|0\rangle$ is annihilated by all soft gluon (and quark) annihilation operators with $\Lambda<k^+<\vee$, and its energy is set to zero. One- and two-gluon states are
\begin{align}
|g^a_i(k)\rangle&=\frac{a^{a\dagger}_i(k)}{(2\pi)^{3/2}}|0\rangle,\qquad |g^a_i(k^+,z)\rangle=\frac{a^{a\dagger}_i(k^+,z)}{(2\pi)^{1/2}}|0\rangle=2\pi\int\frac{d^2\mathbf k}{(2\pi)^2}e^{-i\mathbf k\cdot z}|g^a_i(k)\rangle,\\
|g^a_i(k)\,g^b_j(p)\rangle&=\frac{a^{a\dagger}_i(k)\,a^{b\dagger}_j(p)}{(2\pi)^3}|0\rangle,\qquad |g^a_i(k^+,z)\,g^b_j(p^+,z')\rangle=\frac{a^{a\dagger}_i(k^+,z)\,a^{b\dagger}_j(p^+,z')}{2\pi}|0\rangle,
\end{align}
normalized as $\langle g^b_j(p)|g^a_i(k)\rangle=\delta^{ab}\delta_{ij}\delta^{(3)}(k-p)$ and $\langle g^b_j(p^+,z')|g^a_i(k^+,z)\rangle=\delta^{ab}\delta_{ij}\delta(k^+-p^+)\delta^{(2)}(z-z')$. The free energies are $E_g(k)=\mathbf k^2/2k^+$ and $E_{gg}(k,p)=\mathbf k^2/2k^++\mathbf p^2/2p^+$.

\paragraph*{Matrix elements.} The LCPT matrix elements of the interaction terms of Sec.~\ref{sec:hamiltonian} are:
\begin{itemize}
\item emission of a gluon by the valence charges,
\begin{equation}
\langle g^a_i(k)|H_g|0\rangle=\frac{g\,\mathbf k^i\rho^a(-\mathbf k)}{4\pi^{3/2}|k^+|^{3/2}};
\label{eq:m1}
\end{equation}
\item emission of an extra gluon in the presence of a gluon,
\begin{equation}
\langle g^c_l(q)\,g^b_j(p)|H_g|g^a_i(k)\rangle=\delta^{ac}\delta_{il}\delta^{(3)}(k-q)\frac{g\,\mathbf p^j\rho^b(-\mathbf p)}{4\pi^{3/2}|p^+|^{3/2}}+\delta^{ab}\delta_{ij}\delta^{(3)}(k-p)\frac{g\,\mathbf q^l\rho^c(-\mathbf q)}{4\pi^{3/2}|q^+|^{3/2}};
\label{eq:m2}
\end{equation}
\item the three-gluon vertex,
\begin{align}
\langle g^c_l(q)\,g^b_j(p)|H_{ggg}|g^a_i(k)\rangle={}&\frac{igf^{abc}\delta^{(3)}(k-p-q)}{8\pi^{3/2}\sqrt{k^+p^+q^+}}\Big[\Big(\mathbf p^i-\mathbf q^i+\frac{q^+-p^+}{k^+}\mathbf k^i\Big)\delta_{jl}\notag\\
&\hspace{6em}+\Big(\mathbf k^j+\mathbf q^j-\frac{k^++q^+}{p^+}\mathbf p^j\Big)\delta_{il}\notag\\
&\hspace{6em}+\Big(\frac{k^++p^+}{q^+}\mathbf q^l-\mathbf p^l-\mathbf k^l\Big)\delta_{ij}\Big];
\label{eq:m3}
\end{align}
\item instantaneous emission of two gluons by the valence charges,
\begin{equation}
\langle g^c_j(q)\,g^b_i(p)|H_{gg\text{-inst}}|0\rangle=\frac{ig^2f^{abc}(q^+-p^+)\,\delta_{ij}\,\rho^a(-\mathbf p-\mathbf q)}{2(2\pi)^3\sqrt{p^+q^+}\,(p^++q^+)^2};
\label{eq:m4}
\end{equation}
\item instantaneous interaction of a gluon with the valence charges,
\begin{equation}
\langle g^c_j(q)|H_{gg\text{-inst}}|g^b_i(p)\rangle=\frac{ig^2f^{abc}(p^++q^+)\,\delta_{ij}\,\rho^a(\mathbf p-\mathbf q)}{2(2\pi)^3\sqrt{p^+q^+}\,(p^+-q^+)^2}.
\label{eq:m5}
\end{equation}
\end{itemize}
The valence charges in momentum space are $\rho^a(\mathbf k)=\int_xe^{i\mathbf k\cdot x}\rho^a(x)$, with $[\rho^a(\mathbf k),\rho^b(-\mathbf p)]=if^{abc}\rho^c(\mathbf k-\mathbf p)$.

\paragraph*{Fourier transforms.} The following integrals convert the momentum-space wave functions to coordinate space. Integrals with a pole in the integration region are defined as principal values.
\begin{align}
&\int\frac{d^2\mathbf k}{2\pi}\frac{\mathbf k^i}{\mathbf k^2}e^{i\mathbf k\cdot X}=\frac{iX^i}{X^2},\label{eq:i1}\\
&\int\frac{d^2\mathbf k}{2\pi}\frac{d^2\mathbf p}{2\pi}\frac{\mathbf k^i\mathbf p^j}{\mathbf k^2(a\mathbf k^2+b\mathbf p^2)}e^{i\mathbf k\cdot X+i\mathbf p\cdot Z}=-\frac{X^iZ^j}{Z^2(bX^2+aZ^2)},\label{eq:i2}\\
&\int\frac{d^2\mathbf k}{2\pi}\frac{d^2\mathbf p}{2\pi}\frac{1}{\xi\mathbf k^2-\mathbf p^2}e^{i\mathbf k\cdot X+i\mathbf p\cdot Z}=\frac{1}{\xi Z^2-X^2},\label{eq:i3}\\
&\int\frac{d^2\mathbf k}{2\pi}\frac{d^2\mathbf p}{2\pi}\frac{\mathbf k^m-\mathbf p^m}{(\mathbf k-\mathbf p)^2}\frac{\mathbf p^l-\xi\mathbf k^l}{\mathbf p^2-\xi\mathbf k^2}e^{i\mathbf k\cdot X+i\mathbf p\cdot Y}=-\frac{X^m+\xi Y^m}{X^2-\xi Y^2}\frac{X^l+Y^l}{(X+Y)^2},\label{eq:i4}\\
&\int\frac{d^2\mathbf k}{2\pi}\frac{d^2\mathbf p}{2\pi}\frac{\mathbf k^m-\mathbf p^m}{\mathbf p^2-\xi\mathbf k^2}\frac{\mathbf p^l-\xi\mathbf k^l}{(\mathbf p-\xi\mathbf k)^2}e^{i\mathbf k\cdot X+i\mathbf p\cdot Y}=-\frac{X^l+Y^l}{X^2-\xi Y^2}\frac{X^m+\xi Y^m}{(X+\xi Y)^2},\label{eq:i5}\\
&\int\frac{d^2\mathbf k}{2\pi}\frac{d^2\mathbf p}{2\pi}\frac{1}{\xi(1-\xi)\mathbf k^2+\mathbf p^2}e^{i\mathbf k\cdot X+i\mathbf p\cdot Z}=\frac{1}{X^2+\xi(1-\xi)Z^2},\label{eq:i7}\\
&\int\frac{d^2\mathbf k}{2\pi}\frac{\mathbf k^i}{\mathbf k^2}\log\frac{\mathbf k^2}{\mu^2}\,e^{i\mathbf k\cdot X}=\frac{iX^i}{X^2}\Big(-2\gamma_E-\log\frac{X^2\mu^2}{4}\Big).\label{eq:i9}
\end{align}
Eqs.~\eqref{eq:i3}--\eqref{eq:i5} were checked numerically.

\paragraph*{Colour conventions.} $(T^a)_{bc}=-if^{abc}$, $f^{acd}f^{bcd}=N_c\delta^{ab}$, $(T^eT^e)_{bc}=N_c\delta_{bc}$. The adjoint Wilson lines are real and orthogonal, $U^\dagger=U^T$, and $U^{aa'}U^{bb'}U^{cc'}f^{a'b'c'}=f^{abc}$. A useful identity is $f^{abc}f^{def}U^{be}(z)U^{cf}(z)=N_c\,U^{ad}(z)$.
